\chapter{Présentation de l'Organisme et Contexte du Stage}

\section{Le Groupe OCP : Acteur Mondial des Phosphates}

Fondé en 1920, le \textbf{Groupe OCP} (Office Chérifien des Phosphates) est l'un des leaders mondiaux de l'industrie des engrais et de la valorisation du phosphate. Détenant les plus importantes réserves de phosphate au monde, l'entreprise opère sur l'ensemble de la chaîne de valeur du minerai : de l'extraction minière (gisements de Khouribga, Gantour, Boucraâ) au traitement chimique et à la fabrication d'acide phosphorique et d'engrais spécialisés (plateformes industrielles de Jorf Lasfar et Safi), jusqu'à l'exportation internationale.

\section{Chaîne de Traitement et Processus Industriel}

Le traitement du minerai brut extrait de la mine implique une succession d'opérations mécaniques et physico-chimiques destinées à enrichir sa teneur en $P_2O_5$ (pentoxyde de phosphore) et à éliminer les impuretés stériles (silice, carbonates). 

Le flux standard de production industrielle s'organise en quatre phases clés :
\begin{enumerate}
    \item \textbf{Concassage Primaire \& Secondaire :} Réduction granulométrique des blocs de roche brute ;
    \item \textbf{Criblage \& Séparation Granulométrique :} Tri des fractions selon la taille des particules ;
    \item \textbf{Convoyage Continu :} Acheminement des flux miniers entre les unités de transformation via des réseaux de bandes transporteuses motorisées ;
    \item \textbf{Enrichissement \& Flottation :} Séparation sélective des grains de phosphate par procédé d'agitation et d'aération en milieu aqueux.
\end{enumerate}

Cette chaîne de production présente une dépendance fonctionnelle séquentielle stricte : l'arrêt imprévu d'une machine en amont engendre une rupture d'alimentation matière sur l'ensemble de la ligne aval.

\section{L'Entité d'Accueil et la Stratégie Digitale}

Le stage a été réalisé au sein du département en charge des opérations et de la transformation digitale sur le site industriel de \textbf{[SITE OCP]}. 

Dans le cadre de sa stratégie d'\textit{Excellence Opérationnelle} et de digitalisation de ses actifs industriels, le Groupe OCP encourage le développement de solutions logicielles souveraines à forte valeur ajoutée. L'intégration de prototypes de Jumeaux Numériques vise à doter les équipes d'exploitation d'outils modernes d'aide à la décision, capables de réduire le taux de défaillance non planifiée (\textit{Unplanned Downtime}) et d'accroître le taux de rendement synthétique (TRS) des installations.

\section{Positionnement et Cadre du Projet}

Ce projet s'inscrit dans une démarche de prototypage logiciel avancé (Proof of Concept --- PoC) menée en autonomie. Il a pour but de valider la faisabilité technique, architecturale et algorithmique d'une plateforme de Jumeau Numérique temps réel intégrant l'analyse d'impact topologique, avant son déploiement à grande échelle sur des environnements connectés au bus industriel de l'usine.
